\HMTNarrativeHeading{Proporción, refinamiento y restitución}

La extrema y media razón holográfica reúne aquí la completación
aritmética y la recuperación operatoria. El refinamiento con acarreo
redistribuye una unidad entre componentes y conserva la suma
normalizada. La división euclídea, aplicada con la profundidad y
la convención de acarreo declaradas, recupera las etiquetas
anteriores. La conservación y la inversa reciben así operaciones
concretas.

La norma ternaria aporta el entero setenta y tres y su interpretación
algebraica y reticular. La relación con la completación decimal
se expresa mediante identidades exactas. El balance bilateral
extiende el análisis a dos transportes isométricos entre espacios
de dimensión arbitraria. Los términos cruzados se cancelan en
una combinación ponderada y restituyen la norma cuadrada inicial.

La especialización nonádica determina el reparto inicial de pesos
setenta y dos sobre setenta y tres y uno sobre setenta y tres.
Una mezcla ortogonal produce después los dos canales normalizados.
Cuando los transportes admiten el relativo unitario declarado,
el lector efectivo combina el observable con su conjugado; su
invariancia equivale entonces a la conmutación con ese relativo. El transporte
del estado completo conserva simultáneamente las relaciones
que permite evaluar cada observable.

La política de archivo prolonga este balance. Las memorias
sucesivas y el terminal mantienen la totalidad a profundidad
finita. La estimación uniforme del terminal permite obtener
un archivo infinito recuperable mediante el adjunto del operador
isométrico. Las cotas describen la estabilidad de esa recuperación
cuando aumenta la profundidad.

La constante de acoplamiento participa en la misma construcción.
Su órbita y su forma de Gram permiten identificar operadores;
la isometría de archivo conserva esas respuestas y sus cotas.
El refinamiento de etiquetas se levanta a un transporte unitario
y se compone con el balance bilateral. La unidad conserva así
su contenido aritmético, su realización lineal y la posibilidad
de restitución. Ésta es la culminación operativa del argumento
rector: una estructura generada admite transformaciones que
redistribuyen su lectura y conservan los datos de su recuperación.
